\section{Conclusion}
\label{sec:conclusion}

We introduced GeoSection, a ratio-optimal recursive bisection algorithm that
adds isoperimetric normalisation to the standard minimum-edge-cut approach.
The core innovation --- dividing the edge-cut by $\sqrt{\min(i,k-i)}$ before
comparing ratios --- prevents the caterpillar pathology without eliminating
genuinely compact urban splits.

The 50-state sweep demonstrates two qualitatively different outcomes.
For states with truly compact, geographically isolated metropolitan areas
(Wisconsin: Milwaukee, Pennsylvania: Philadelphia, Tennessee: Nashville/Memphis),
the normalisation \emph{confirms} the asymmetric split: these urban cores
are genuinely compact by the isoperimetric standard.
For states with distributed population (North Carolina, Maryland), the
normalisation selects a near-equal split: the urban cores are not compact
enough to dominate the isoperimetric comparison.

\textbf{The partisan outcome is robust across bisection strategies.}
Of 44 comparable states, 38 produce identical seat counts under GeoSection
and the B.7 MEC baseline.
The mean proportionality gap in competitive states is $-7.0$\,pp under
GeoSection, indistinguishable from the B.7 baseline.
Geographic voter sorting --- not the bisection algorithm --- is the primary
determinant of partisan outcomes.
This is GeoSection's strongest empirical result: changing the bisection
ratio structure does not change the number of Democratic seats in 86\,\%
of states.

\textbf{GeoSection is legally defensible.}
The natural ratio is computable from public data (TIGER boundaries, census
populations) via a single published formula.
It does not require partisan data, and the result can be independently
verified by any party.
States where the peel is genuine have a stronger legal defence than under
standard bisection: ``geography confirmed this ratio via isoperimetric
normalisation'' is a more principled claim than ``this ratio minimises
edge-cut.''

\paragraph{Future work.}
Three directions merit investigation:

\begin{itemize}
\item \textbf{Phase 2 evaluation.}
  The directional penalty $\lambda$ is implemented but not characterised.
  Systematic tests across the 50 states would quantify the tradeoff between
  boundary straightness and edge-cut cost.

\item \textbf{Cross-census stability.}
  Computing natural ratios for 2000 and 2010 Census data would test whether
  the geographic structure of states is stable across apportionment cycles,
  or whether urbanisation trends shift the natural ratio over time.

\item \textbf{Higher seed counts for large states.}
  California ($k=52$), Texas ($k=38$), and Florida ($k=28$) have uncertain
  natural ratios at 50 seeds per ratio.
  Running 200 seeds per ratio would improve confidence in the convergence
  of the selected ratio for these large states.
\end{itemize}

\paragraph{Code and data availability.}
GeoSection is implemented in the open-source \texttt{redist} Rust binary
(\texttt{redist state --partition-mode geosection}).
50-state tract-level assignments and natural ratio logs are available in the
replication archive (\texttt{research/A.4+replication-materials}).
